%We first present a brief description of the fundamental SEIR model to provide the epidemiological context of our work. Next, we discuss related research, mainly in the form of modeling tools that have been used in the research community of disease modeling.



\subsection{Model-Driven Engineering and Domain-Specific Languages}

Model-driven engineering (MDE) has seen significant adoption across various domains beyond epidemiology, demonstrating its versatility and effectiveness. Recent systematic reviews have examined the intersection of MDE with emerging technologies and application domains. Rädler et al.~\cite{radler2024bridgingmdeai} conducted a systematic review of domain-specific languages and model-driven practices in AI software systems engineering, identifying how MDE techniques can tame the complexity of implementing AI algorithms. Similarly, Naveed et al.~\cite{naveed2024mde4ml} performed a systematic literature review on model-driven engineering for machine learning components, highlighting the benefits of MDE in managing the complexity of ML-enabled systems. Zafar et al.~\cite{zafar2024dsmmde} explored the effectiveness and trends of domain-specific model-driven engineering through a systematic literature review, emphasizing benefits such as enhanced efficiency, improved maintenance, and reduced development time.

The application of MDE to healthcare and medical domains has also gained traction. Trevena et al.~\cite{trevena2024digitaltwins} presented a model-driven approach for digital twins in patient simulation, using graph-based models to create virtual patient replicas for testing clinical interventions. Brandon et al.~\cite{brandon2023mdehealth} reviewed model-driven development for AI-based healthcare systems, demonstrating applications ranging from disease diagnosis to patient information management. These works share similarities with our epidemiological modeling approach, particularly in managing complex domain-specific knowledge through metamodeling.

Domain-specific languages (DSLs) for data-intensive workflows have also evolved. Salado-Cid et al.~\cite{saladocid2023swel} introduced SWEL, a platform-independent DSL for modeling data-intensive scientific workflows, addressing interoperability challenges similar to those faced in epidemiological modeling. The recent integration of large language models with MDE has opened new possibilities. Di Rocco et al.~\cite{dirocco2025llmmde} examined how LLMs can automate MDE tasks such as model classification and recommendation, potentially lowering barriers to adoption.

An important consideration in MDE adoption is the human factor. Liebel et al.~\cite{liebel2024humanfactors} emphasized that while MDE aims to improve communication and understandability through models, human factors play a critical role in the success of modeling efforts. This aligns with our motivation for EpiMDELite, which prioritizes immediate usability for epidemiologists by using familiar domain concepts rather than requiring extensive MDE expertise.

These advances in MDE demonstrate its growing maturity and applicability across domains. However, epidemiological modeling presents unique challenges related to population stratification, disease dynamics, and the need for rapid model adaptation during evolving epidemics. Our work applies MDE principles specifically to address these epidemiological challenges, providing both an accessible entry point (EpiMDELite) and advanced extensibility for sophisticated modeling needs.

\subsection{Modeling tools}

Spatio-Temporal Epidemiological Modeler \footnote{\url{https://projects.eclipse.org/projects/technology.stem}}, STEM \cite{STEM}, is part of the Eclipse ecosystem and the only epidemiology tool based on Eclipse, along with ours. As the name suggests, it is used to model the spread of disease in a particular geographic region over a given period. The tool requires significant knowledge and assumptions from the perspective of the modeller. Its interface is not always intuitive, which can lead to errors that cause the platform to crash. This is possibly due to the fact that it integrates many features including modeling, scenario design, parameterization and simulation, which is an overload when trying to model simple cases. Creating a model from scratch is an expensive task and it is easier to adapt existing models. It is a tool specifically tailored for disease modeling with geographic considerations. Nevertheless, STEM is backed by a large community and it is good for large scale scenarios, such as the Global Air Travel Network~\cite{STEM}, which models passenger travel through commercial airports.

The Institute for Disease Modeling Compartmental Modeling Software, IDM-CMS, allows users to run simple compartmental models. Each compartment has to be defined individually, so there is no option for Cartesian product of graphs. The software offers no graphical interface, it is a command line tool, making it error-prone, since the input files take hundreds of lines that have to be written by hand \footnote{\url{https://github.com/InstituteforDiseaseModeling/IDM-CMS/blob/master/models/idm}}. Our consulting expert has confirmed that this the usual case for most tools. He explained that for custom models, a modeler uses tools like Matlab and has to write a great amount of similar lines by hand, which involves copy-pasting. Naturally, if the model is not stratified, there would be minimal copy-pasting, but one of our use cases is models with a large amount of stratifications, so this software is not directly applicable. The input format for the equations is lisp-style, as is in the case of our platform, so in principle our platform can support models created in the IDM-CMS tool. In addition, there is no notion of configuration in IDM-CMS, but an already parameterized model is required, while in EpiMDE, we decouple the model from its parameterization. This can improve the reusability of the model in multiple simulation scenarios.

Climate-Sensitive Infectious Disease, CSID, often refers to diseases which can spread to humans from insects, food or water during certain periods of the year depending on meteorological conditions. A study by the Wellcome Institute\footnote{\url{https://wellcome.org/}} identifies 37 tools for CSID modeling \cite{Wellcome2022}. The study only identifies the tools and characterizes them based on the information available without using the tool. The most interesting project identified in the report is ANOSPEX: a stochastic, spatially explicit model for studying Anopheles meta-population dynamics. It is a model of the Anopheles marsh mosquito populations based on climate conditions. While this tool is very specific within the phenomenon it studies, these types of models are in principle easy to construct using EpiMDE, our proposed platform.

Kendrick \cite{Kendrick} is a DSL and simulation platform for epidemiology modeling\footnote{\url{https://github.com/KendrickOrg/kendrick}}. While this description is an exact match for EpiMDE, Kendrick's DSL is textual, not graphical. There are other differences from our platform. Firstly, structure, equations, parameters and simulation details are in the same file.
Secondly, stratification seems possible according to the second example model\footnote{\href{https://github.com/KendrickOrg/kendrick/blob/c84ad7ba63b42332efed1ef66aaceba0566d2549/documentation/formal-SoC-models/m2.md?plain=1\#L24-L26}{Kendrick formal-SoC-model m2}} but it is not as clear\footnote{\href{https://github.com/KendrickOrg/kendrick/blob/c84ad7ba63b42332efed1ef66aaceba0566d2549/documentation/meta-model/meta-modelv3.org?plain=1\#L47-L52}{Kendrick MetaModel}}.
Finally, the DSL may be a form of MDE, but there is no explicit use of a MDE framework.
The project is implemented in Small-Talk and runs in Pharo\footnote{\url{https://pharo.org/}} using Metacello\footnote{\url{https://github.com/Metacello/metacello}}.
Much like many projects, it is hard to verify the features presented as the instructions to run are incorrect and there is no output from Pharo.

Our EpiMDE platform is distinct from the EpiModel R package~\cite{jenness2018epimodel}. While both tools aim to facilitate epidemiological modeling, they differ significantly in their approach and implementation. The EpiMDE platform is based on model-driven engineering principles and technologies, such as the Eclipse Modeling Framework (EMF) and Sirius for creating and managing models, offering a graphical interface and supporting extensibility through metamodeling. In contrast, the EpiModel R package is a mathematical and statistical tool designed for epidemiological analysis and simulation within the R ecosystem, focusing on providing mathematical functions for compartmental and stochastic individual contact models and data analysis capabilities. There is no direct technical relationship between our platform and the R package, although both serve complementary purposes in the field of epidemiology.

More recently, Lemaitre et al.~\cite{lemaitre2024flepimop} presented \emph{flemiMoP}, a flexible and open-source pipeline to support compartmental modeling for COVID-19. The pipeline is modular, consisting of modules in R and python, and the users can specify different models through a YAML configuration file. The tool focuses almost exclusively on the simulation component of modeling, also including parameter inferencing. Compared to EpiMDE, flemiMoP still requires high epidemiological expertise, it does not include any graphical representation of the models to facilitate communication, while if there is a need to modify a module, good understanding of R or python is required.

\subsection{Large language models in epidemiological modeling}

In epidemiology, LLM applications extend beyond traditional modeling to disease surveillance, outbreak detection, and real-time forecasting~\cite{b12}. Notable developments include PandemicLLM~\cite{b14}, which uses multi-modal LLMs for disease forecasting by incorporating complex non-numerical information, and infodemiology applications for epidemic assessment from social media~\cite{b15,b16}. The integration of AI with mechanistic epidemiological modeling represents a transformative approach, combining traditional epidemiological knowledge with modern data-mining capabilities to enhance pandemic preparedness~\cite{b17}. The majority of current LLM applications in epidemiology focuses on forecasting and as such it differs from EpiMDE's use of LLMs in two important ways: a) other works require existing transmission data as input, while EpiMDE requires only structural and parametric data as input, focusing mostly on the modeling part, and b) other works almost entirely bypass the modeling component, and its graphical representation, which significantly deprives researchers from the ability to share, communicate and macroscopically investigate the structure of the model, and more importantly the relation between the compartments. In practice, EpiMDE's LLM component is used one stage earlier than other LLMs and in this way, it can work complementary to other works.

